\documentclass[../main]{subfiles}
\begin{document}
\section{Exposición y Documentación del Trabajo Integrador de
Máquinas Eléctricas}

\info{Enfoque Específico y Rigor Técnico}{El presente trabajo tiene
  como único y exclusivo objeto de estudio la sección correspondiente a
  las máquinas eléctricas dentro del desarrollo integral de los alumnos
  de 6to año. No se evalúan aspectos ajenos al electromagnetismo,
  automatismos genéricos o estructuras mecánicas secundarias; el rigor
  técnico se centra estrictamente en el dominio de la conversión
electromecánica de energía y sus leyes fundamentales.}

\subsection{1. Guía de Identificación de Máquinas Eléctricas en el Proyecto}
Para dar inicio al desarrollo, cada grupo debe realizar un riguroso
relevamiento y justificación de las máquinas eléctricas presentes o
necesarias en su diseño. A modo de guía orientativa, todo sistema
técnico de 6to año involucra al menos una de las siguientes
categorías tecnológicas:
\begin{itemize}
  \item Motores de Corriente Alterna (Asincrónicos y Síncronos):
    Utilizados en tracción, bombeo, transporte o accionamientos
    generales, abarcando motores de jaula de ardilla, rotor bobinado
    y motores monofásicos de fase partida.
  \item Motores de Corriente Continua y Accionamientos Especiales:
    Servomotores, motores paso a paso (steppers), motores brushless
    (sin escobillas) o motores universales aplicados en
    posicionamiento de alta precisión y control dinámico.
  \item Transformadores y Elementos Estáticos: Unidades destinadas a
    la transformación de niveles de tensión, aislamiento galvánico,
    adaptación de impedancias o conversión estática de energía.
  \item Generadores y Sistemas de Recuperación: Máquinas rotativas
    destinadas a la generación de energía eléctrica o sistemas de
    frenado regenerativo e inercial.
\end{itemize}

\subsection{2. Estructura de la Documentación Técnica}
La presentación formal de la documentación técnica (desterrando la
denominación informal de informe) debe estructurarse rigurosamente
bajo los siguientes apartados normativos:
\begin{enumerate}
  \item Carátula y Datos Institucionales: Identificación precisa del
    equipo de trabajo, división, especialidad, ciclo lectivo y
    denominación exacta del sistema electromecánico.
  \item Identificación y Justificación Funcional: Descripción
    detallada del rol que cumple cada máquina eléctrica en el
    proyecto, justificando su elección frente a otras alternativas
    tecnológicas en función de par, velocidad y rendimiento.
  \item Memoria de Cálculo Electromecánico:
    \begin{itemize}
      \item Parámetros constructivos y eléctricos: tensiones
        nominales, corrientes de servicio y arranque, frecuencias,
        velocidad sincrónica y deslizamiento.
      \item Análisis energético integral: cálculo de potencias,
        factor de potencia y rendimiento global.
      \item Dimensionamiento y selección normada de conductores,
        protecciones termomagnéticas, diferenciales y dispositivos de
        maniobra (contactores, relés térmicos).
    \end{itemize}
  \item Esquemas y Diagramas Eléctricos: Planos constructivos en CAD,
    diagramas unifilares y circuitos de fuerza y comando completos,
    normalizados y debidamente acotados.
  \item Conclusiones y Normativa Aplicada: Evaluación crítica del
    comportamiento operativo del sistema frente a las exigencias
    reales y cumplimiento estricto de normativas vigentes (IRAM / IEC).
\end{enumerate}

\subsection{3. Pautas y Recomendaciones para la Exposición Oral}
\begin{itemize}
  \item Enfoque Exclusivo en Máquinas: La defensa oral no debe
    relatar la historia del proyecto general, sino concentrarse de
    forma quirúrgica en el sistema de máquinas eléctricas, sus
    ecuaciones, su dimensionamiento y su comportamiento dinámico.
  \item Soporte Visual Minimalista: Las diapositivas deben contener
    exclusivamente esquemas eléctricos, diagramas vectoriales,
    circuitos equivalentes y ecuaciones fundamentales. Está
    terminantemente prohibido leer texto corrido.
  \item Previsión Logística: Solicitar los recursos audiovisuales
    (TV/Proyector) con la debida anticipación y llevar el material
    respaldado en un pendrive o medio local para prevenir fallas de
    conectividad.
  \item Gestión Temporal y Dominio Técnico: Cada equipo dispondrá de
    un tiempo estricto de 15 a 20 minutos. Todos los integrantes
    deben demostrar idoneidad absoluta para responder
    cuestionamientos técnicos sobre cualquier cálculo o selección realizada.
\end{itemize}
\end{document}
